\chapter{力的特性}

\section{什么是力}

物理定律能够帮助我们认识自然与利用自然。虽然仅仅从这一点看，就已经值得我们花时间去研究它，但是人们还是应当不时停下来思考一下“它们的真正含义是什么？”自远古以来，任何领域所表述的含义都是一个使哲学家发生兴趣而又感到棘手的课题，而物理定律的含义甚至更为有趣，这是因为人们普遍认为，这些定律阐明了某种真正的知识。知识的含义是一个深奥的哲学问题，而追问“它的含义是什么？”始终是重要的。

让我们来问一下，“写成$F=ma$的牛顿定律的含义是什么？力、质量以及加速度的含义是什么？”嗯，我们可以凭直觉领会质量的含义，如果我们知道了位置和时间的含义，就能够定义加速度。我们将不去讨论这些含义，而要集中讨论力这个新概念。答案同样地简单：“如果一个物体在作加速运动，那么一定有力作用在这个物体上。”这就是牛顿定律所表明的，于是人们可以想象到的力的最精确和最优美的定义就能够简单地表述为：力等于物体的质量乘以加速度。假定我们有一条定律，即：如果所有外力之和等于零，那么动量守恒成立；于是产生了这样的问题，“所有外力之和等于零是什么意思？”对上面这个陈述的一个有意思的解释是：“当总动量不变时，外力之和等于零。”这种说法肯定有毛病，因为它实在没有说出什么新的东西。如果我们已经发现了一条基本定律，该定律宣称：力等于质量乘以加速度，并且正式加以定义，那么我们就什么也没有发现。我们也可以这样来定义力：一个物体不受力作用时将一直以恒定速度沿直线运动。于是，如果我们观察到一个物体不是以恒定速度沿直线运动，我们就可以说，物体上有力作用着。这些说法当然不可能属于物理学的内容，因为它们是一些循环论证的定义。上面的牛顿表述尽管看来好像是力的一种最精确的定义，而且合数学家的意；然而，它完全是无用的，因为从这个定义得不出任何预见。人们可以整天坐在安乐椅上，任意定义一些词，但是要发现当两个球相互碰撞或者一个重物悬挂在弹簧上时发生一些什么则完全是另一回事，因为物体所表现的行为是与定义的选择完全无关的。

举例来说，如果我们愿意说，一个物体不去碰它，它就停留在原来所处的位置上不动，那么当我们看到某个东西在漂移时，我们就可以说，这一定是由一种称为“戈斯”(Gorce)的力\footnote{此名称系作者根据“Force(力)”杜撰}所引起的——“戈斯”是位置的变化率。现在我们得到了一个奇妙的新的定律：任何物体都保持静止状态除非有“戈斯”作用其上。你们看，这就和上述力的定义相类似，它并不包含什么内容。牛顿定律的真正内容是：除了$F=ma$这条定律外，力还应该具有某些独立的性质。但是牛顿或其他人并没有把力所特有的独立性质完全描述出来。因此物理定律$F=ma$是一个不完全的定律。它暗示着，如果我们研究质量乘以加速度，并且称这个乘积为力，也就是说如果我们有兴趣研究力的特征，那么我们将发现力具有某种简单性。这个定律是分析自然界的一种很好的方案，力是简单的则是一种联想。

这类力的第一个例子就是完整的万有引力定律，它是由牛顿提出的，并且在阐述定律时还回答了“力是什么？”这个问题。假如除了万有引力外没有其他的力，那么万有引力定律和力的定律（第二运动定律）的组合应当是一种完整的理论。但是除了万有引力外还有很多种力，而且我们希望在许多不同的情况下使用牛顿定律，因此为了继续深入，我们必须讲一些有关力的性质的内容。

例如，在处理力的问题时总是默认：除非有某个物理实体存在，否则力等于零。如果我们发现一个力不等于零，我们也一定能在附近找到某个物体是力的来源。这个假定完全不同于我们上面所介绍的“戈斯”的那种情况。力的最重要的特征之一就是它具有实质性的起源，这就不仅是一个定义了。

牛顿还提供了一个有关力的法则：相互作用的物体之间的力大小相等、方向相反——作用力等于反作用力。后来知道，这条法则并不是完全正确的。事实上，$F=ma$这条定理就不完全正确。要是它是一个定义，我们就应当说：它永远是完全正确的，但事实上并非如此。读者可能会提出异议：“我不喜欢这种不严谨性，我喜欢对每件事情都下一个严格的定义；事实上在有些书上是这样说的，任何科学都是一门严格的学科，在那里每件事情都有一个定义。”如果你坚持要获得力的精确定义，你将永远得不到它！因为首先，牛顿第二定律是不精确的，其次，要理解物理定律，你就必须懂得所有这些物理定律都是某种近似。

读者可能会提出异议：“我不喜欢这种不严谨性，我喜欢对每件事情都下一个严格的定义；事实上在有些书上是这样说的，任何科学都是一门严格的学科，在那里每件事情都有一个定义。”如果你坚持要获得力的精确定义，你将永远得不到它！因为首先，牛顿第二定律是不精确的，其次，要理解物理定律，你就必须懂得所有这些物理定律都是某种近似。任何简单的概念都是近似的。作为例子，考虑一个客体，什么是客体？哲学家总是这样说：“嗯，就拿一张椅子来作为例子吧！”当他们说这句话的时候，你们就知道他们不知道再如何说下去了。椅子是什么？椅子就是那边的某种确定的东西，确定到什么程度呢？原子不时从椅子中跑出来——虽然不是很多，但总有一些，灰尘落到椅子上，逐渐溶化到油漆里，所以要精确地定义椅子，确切地说出哪些原子属于椅子，哪些原子属于空气或哪些原子属于灰尘，哪些原子属于椅子上的油漆，这是不可能的。因此只能近似地定义椅子的质量。同样，要确定单个客体的质量也是不可能的。因为世界上并不存在任何单个孤立的客体——每个客体都是许多事物的混合体。所以我们只能以一系列的近似和理想化来处理它们。

任何简单的概念都是近似的。作为例子，考虑一个客体，什么是客体？哲学家总是这样说：“嗯，就拿一张椅子来作为例子吧！”当他们说这句话的时候，你们就知道他们不知道再如何说下去了。椅子是什么？椅子就是那边的某种确定的东西，确定到什么程度呢？原子不时从椅子中跑出来——虽然不是很多，但总有一些，灰尘落到椅子上，逐渐溶化到油漆里，所以要精确地定义椅子，确切地说出哪些原子属于椅子，哪些原子属于空气或哪些原子属于灰尘，哪些原子属于椅子上的油漆，这是不可能的。因此只能近似地定义椅子的质量。同样，要确定单个客体的质量也是不可能的。因为世界上并不存在任何单个孤立的客体——每个客体都是许多事物的混合体。所以我们只能以一系列的近似和理想化来处理它们。

窍门就在于使之理想化。在大约是$\frac{1}{10^{10}}$以内的极好近似下，椅子中的原子数在一分钟内是不变的，而如果要求不是太精确的话，我们可以把椅子理想化为一个确定的物体；同样，如果要求不是太精确的话，我们将以一种理想化的方式来学习有关力的特性。人们可能并不满足于物理学试图获得的对自然界的近似（但总是力图增加这个近似的精确度）看法，而宁愿要一种数学的定义，但是数学的定义在真实世界中是永远不会起作用的。对数学来说所有的逻辑都能完全贯彻，数学定义对它将是适合的。但是正如我们在像海浪和葡萄酒等一些例子中指出过的那样，物理世界是复杂的。当我们试图将它的各部分隔离开来讨论一种质量时，对于酒和杯子来说，当一方溶于另一方时，我们怎么能知道哪些是酒，哪些是杯子呢？作用于单个物体上的力已经包含着近似了，要是我们有一个论述真实世界的体系，那么这种体系，至少对目前来说，必定包含着某些近似。

这种体系完全不像数学体系那样。在数学中对每个事物都能够下定义，此外我们就不知道在说些什么了。事实上，数学的伟大就在于我们不必说出我们所谈论的是什么东西。此种伟大在于定律、论证和逻辑都与“它”是什么东西无关。如果有另外一组客体，它遵从同样的欧几里得几何公理体系，那么一旦我们作出一些新的定义，并根据它们作正确的逻辑推演，所有的结论都将是正确的，而与所讨论的内容没有关系。然而，在自然界中当我们利用一束光或经纬仪（正如我们在测量中所做的）来画出或建立一条线时，我们是在欧几里得的几何意义上量度线吗？不，我们是在作一种近似。叉丝有一定的宽度，但是一根几何线并没有宽度。因此，欧几里得几何学是否能用于测量是一个物理问题而不是数学问题。然而，从实验的观点，而不是从数学的观点来看问题的话，我们需要知道欧几里得的定理是否适用于我们丈量土地时所使用的那种几何。于是我们假定它是适用的，而且非常成功。但是它不是精确的，因为我们测量的线并不是真正的几何线。那些真正抽象的欧几里得线是否适用于实验上的线是一个经验问题；而不是由纯粹的推理可以回答的问题。

同样，既然力学是一门描述自然界的学科，我们就不能仅把$F=ma$称作定义，纯数学地推演一切，把力学变成一门数学理论。在建立了适当的假定后，总是有可能建立一套数学体系，就像欧几里得所作的那样，但是我们不可能建立有关我们这个世界的某种数学，因为迟早我们必定会发现这些公理对于自然界的客体是否正确。因此我们立刻被纠缠入这些复杂的、“玷污了的”自然界的客体中去，而且作出精确度日益增加的近似。

\section{摩擦力}

前面的考虑表明，要真正领会牛顿定律需要对力进行讨论。本章的目的正是要提出这种讨论使牛顿定律更为完整。我们已经研究了加速度的定义和有关的概念，但是现在我们必须研究力的性质。与前面几章不同，这一章不是非常严谨的，因为力是十分复杂的。

首先从特殊的力开始，我们来考虑作用在空中飞行的飞机上的阻力。这种力遵从什么定律（的确，对每一种力都有一条定律，也一定得有一个定律）？人们简直想不到这种力所遵从的定律会如此简单。试想一想作用在空中飞行的飞机上的阻力是由什么构成的——冲过机翼的空气，机尾的漩涡，机身周围发生的变化以及许多其他复杂因素。于是你们就认为将不会有一个简单的定律。但是作用在飞机上的阻力近似地等于一个常数乘以速度的平方，或$F\approx cv^2$，则是一个显著的事实。

那么这样一个定律的地位如何呢？它类似于$F=ma$吗？绝对不是。首先，因为这一定律是由风洞试验大致得出的经验公式。你也许会说：“好！$F=ma$也可能是经验的。”那不是两者有所差别的原因。差别不在于它是经验的，而在于就我们所了解的自然界来说，这个定律是事件的错综复杂性的产物，它根本不简单。如果我们研究得越来越深入，测量得越来越精确，这个定律就变得越加复杂，而不是越加简单。换句话说，当我们越来越细致地研究这个飞机阻力的定律时，我们发现它越来越“不真实”，而我们研究得越深入，测量得越精确，真相就变得越加复杂。因此在这种意义上，我们认为它不是由那种简单的基本过程所引起的，与我们原来的推测一致。举例来说，如果速度非常低，低到一般的飞机不能飞行，那么当飞机在空气中被慢慢地拖着前进时，定律就发生变化，这时曳引阻力与速度之间的关系比较接近于线性。再举另一个例子，球、气泡或任意物体在蜂蜜那样的黏稠液体中缓慢地运动时，作用于其上的摩擦阻力同速度成正比。但是当运动速度变快，以致引起液体打漩时（蜂蜜不会打漩，但水和空气会打漩），那么摩擦阻力就更接近于同速度的平方成正比（$F=cv^2$），而如果速度继续增大，甚至这个定律也开始失效。有人说：“这是由于比例系数有了某些改变”，说这些话的人是回避问题。其次，还有其他很复杂之处：作用在飞机上的力是否可以分成或分解为一个作用在机翼上的力，一个作用在机头上的力等等？的确，如果我们考虑的是作用在这里或那里的力矩的话，是可以这样做的。但这时我们必须找出作用在机翼上等等的力的特殊定律。令人惊异的是，作用于一个机翼上的力与作用于另一个机翼上的力有关。换句话说，如果我们把飞机拆开，只把一个机翼放到空中，这时机翼所受的力与飞机的其余部分都在那里时所受的力是不同的。当然，这是因为有一些冲击机头的风流到机翼周围，改变了作用在机翼上的力。这样一种简单的、粗糙的、经验的定律能够应用于飞机设计中似乎是一个奇迹，但是这个定律与物理学的基本定律不是同一类型的，进一步的研究只会使这种定律越来越复杂。对于系数$c$如何依赖于飞机前缘形状的研究，说得婉转一点，是令人失望的。根据飞机形状来决定系数的简单定律是不存在的。对比之下，万有引力定律是简单的，越是研究就越觉得它是简单的。

刚才我们讨论了物体在空气中快速运动和在蜂蜜里缓慢运动而引起的两种摩擦情况。另外还有一类摩擦，它是当一个固体在另一个固体上滑动时出现的，称为干摩擦或滑动摩擦。在这种情况下，需要有力来维持运动。这种力称为摩擦力，它的起因也是非常复杂的问题。从原子情况来看，相互接触的两个表面是不平整的。它们有许多接触点，在这些接触点上，原子好像黏结在一起，于是当我们拉动一个正在滑动的物体时，原子啪的一下分开，随即发生振动；所发生的情况大致如此。过去，把这种摩擦的机理想象得非常简单：表面只不过布满凹凸不平的形状，摩擦起因于抬高滑动体越过突起部分。但是不可能是这样，因为在这种过程中不会有能量损失，而实际上是要消耗动力的。动力耗损的机理是当滑动体撞击突起部分时，突起部分发生形变，接着在两个物体中产生波和原子运动，过了一会儿，产生热。现在非常出乎意外的是，根据经验这个摩擦再次可以近似地用一个简单的定律来描述。这个定律是：克服摩擦，使一个物体在另一物体上运动所需的力取决于两个相互接触的表面间的法向力（即同表面垂直的力）。实际上，作为相当好的近似，摩擦力与这个法向力成正比，比例系数近似是常数，即
\begin{equation}
    \label{Eq:I:12:1}
    F=\mu N
\end{equation}
\begin{wrapfigure}{r}{0.4\textwidth}
    \centering
    \includegraphics[width=0.35\textwidth]{Chapter12/接触面相对光滑时的摩擦力与法向力的关系}
    \caption{接触面相对光滑时的摩擦力与法向力的关系}
    \label{figure:接触面相对光滑时的摩擦力与法向力的关系}
\end{wrapfigure}
式中的$\mu$称为摩擦系数（图12-1）。

虽然这个系数不是严格的常数，但是这个公式对于大致判断某些实际或工程学情况中所需力的大小是一个很好的经验法则。如果法向力或运动速度太大，由于产生大量的热，定律就失效。重要的是要认识到每个经验定律都有它的适用范围，超过了这个范围，定律实际上不起作用。

公式$F=\mu N$是近似正确的，这可以用一个简单的实验来验证。取一块平板，使它倾斜一个小角度$\theta$，在平板上放置一块重为$W$的木块，然后使平板倾角增大，直到木块由于自身的重量刚好开始滑动。重力沿平板向下的分力是$W sin \theta$，当木块匀速滑动时，这个分力必须等于摩擦力$F$。与平板垂直的分力是$W cos \theta$，这个力就是法向力$N$。代入这些值后，公式变成$W sin \theta=\mu W cos \theta$，由此我们得到$\mu =\frac{sin \theta}{cos \theta}=tan \theta$。如果这个定律是绝对正确的，在某个一定的倾角，物体就开始下滑。如果在同一木块上再增加一些重量，这时虽然$W$增大，但是公式中的各个力都按相同的比例增加，$W$被消去了。如果$\mu$保持不变，则加重的木块将再次在同样的斜度时下滑。当我们用原先的重物做试验，确定角度$\theta$时，可以发现，比较重的木块将以大致相同的角度下滑，甚至当一个重物比另一重物大很多倍时，上述规律仍然正确。于是我们得出结论：摩擦系数与重量无关。

在做这个实验时，值得注意的是，当平板倾斜到约为正确的角度$\theta$时，木块并不平稳地滑动，而是以断断续续的方式下滑。在某个地方它可能停住，在另一个地方它可能作加速运动。这个现象表明，摩擦系数仅仅大致上是常数，在平板上不同的地点摩擦系数是不同的。不论木块上加重物或不加重物都观察到同样的古怪现象。这些变化是由平板的不同平滑度或硬度引起的，也可能是污物、氧化物或其他外来物质引起的。表中列出的所谓钢与钢、铜与铜等的$\mu$值全都是不可信的，因为它们忽略了上述那些真正决定$\mu$值的因素。摩擦决不是由于“铜与铜”等等引起的，而是由于黏附到铜上的杂质所引起的。

在上述这类实验中，摩擦几乎与速度无关。很多人认为，使物体起动所需克服的摩擦力（静摩擦）大于保持物体滑动所需的力（动摩擦）。但是用干燥金属很难显示出有什么差别。这种见解可能是由于这样的经验引起的：存在着极少量油或润滑剂，或者木块被弹簧或其他易形变的支座撑起以至于看起来好像是结合在一起的。

尽管有大量精确分析的工程数据，要做精确的定量摩擦实验仍然是非常困难的，并且摩擦定律也还没有被人们很好地分析过。虽然，只要表面经受标准处理，定律$F=\mu N$是相当精确的，但是对于定律具有这种形式的原因还没有真正弄清楚。要证明摩擦系数$\mu$几乎与速度无关需要一些巧妙的实验方法，因为如果下表面快速振动，表面摩擦就要大大减少。当在非常高的速度下进行这项实验时，必须注意物体彼此之间不能有振动，因为在高速情形下摩擦明显减小常常是由振动引起的。无论如何，摩擦定律是又一个半经验定律，这些半经验定律还没有被我们完全认识。而且，奇怪的是从我们做过的所有研究来看，对这些现象仍然没有进一步理解。实际上，目前甚至要估计一下两个物体之间的摩擦系数也是不可能的。

上面曾经指出，企图用纯的物体如铜在铜上面的滑动来测量$\mu$将得出虚假的结果。因为接触的表面不单纯是铜，而是氧化物和其他杂质的混合物。如果我们想要得到绝对纯的铜，即使清洗和抛光表面，在真空中对材料除气，并且采取各种可能想到的预防措施，我们还是不能测得$\mu$。因为即使我们把装置倾斜到垂直位置，滑块仍不下落——两片铜粘在一起了！对一般硬度的表面来说，摩擦系数$\mu$通常比1小，而这时$\mu$变得比1大上好几倍！出现这种意想不到的现象的原因是当相互接触的原子全都是同一种原子时，这些原子无法“知道”它们是在不同的铜片上的。当原子存在于其他氧化物、油脂和更复杂的玷污物的薄表面层时，原子就“知道”它们不是在同一部分。当我们考虑到正是原子之间的力把铜原子结合在一起成为固体的时候，就会明白，对纯金属是不可能得出正确的摩擦系数的。

在用一块平玻璃板和一只玻璃杯做的简单的家庭实验中也能够观察到同样的现象。如果把玻璃杯放在平板上，并用一根绳子拉它，它将在平板上很好地滑动，人们能够感觉到摩擦系数，这个系数虽然有点不规则，但毕竟是一个系数。如果我们现在把玻璃板和玻璃杯底弄湿，重新再拉，就会发现杯子粘住了。如果仔细看一下，将会发现划痕，因为水能够从表面除去油脂和其他玷污物，这样我们才真正得到玻璃与玻璃的接触。这种接触是如此牢固，以至于使它们黏紧在一起，不再分离，结果玻璃被撕裂，也就是说造成了划痕。

\section{分子力}

接下来我们将讨论分子力的特征。这些力是原子之间的力，也是摩擦的根本起因。在
\begin{wrapfigure}{r}{0.4\textwidth}
    \centering
    \includegraphics[width=0.35\textwidth]{Chapter12/两个原子之间的力与其距离的函数关系}
    \caption{两个原子之间的力与其距离的函数关系}
    \label{figure:两个原子之间的力与其距离的函数关系}
\end{wrapfigure}
经典物理学的基础上，分子力从来没有得到满意的解释；只有用量子力学才能充分理解它们。然而，根据经验，原子之间的力可用图12.2来说明。

图中将两个原子之间的力$F$作为两个原子之间的距离$r$的函数。同时，还存在着不同的情况：例如在水分子中氧带有较多负电荷，所以负电荷和正电荷的平均位置不在同一点上，结果附近的另一个分子感受到比较大的力，这个力称为偶极-偶极力。然而，对许多系统来说，电荷平衡得非常好，特别是氧气，它是完全对称的。在这种情况下，虽然负电荷和正电荷散布在整个分子中，但是这种分布使正、负电荷的中心重合。正、负电荷中心不重合的分子称为极性分子，电荷与电荷中心间距离的乘积称为偶极矩。非极性分子是一种电荷中心重合的分子。对于所有非极性分子（其中所有的电力都被中和），在较大距离上的作用力仍然是引力，而且与距离的7次方成反比，即$F=\frac{k}{r^7}$，式中$k$是一个取决于分子的常数。只有当我们学到量子力学时才会懂得为什么是这样的。当有偶极子存在时，力就增大。当原子或分子靠得太近时，它们以很大的斥力相互排斥；正是这个力使得我们不会落到地板下面去！

这些分子力可以用一种相当直接的方式来演示：用一只滑动的玻璃杯做的摩擦实验就是方法之一，另一种方法是取两个经过非常仔细研磨和十分平整的平面，使之紧密地贴合在一起。约翰逊(Johansson)平板就是这种平面的一个实例。机床厂里常用它作为精确的长度测量的标准。如果一块这样的平板在另一块平板上非常小心地滑动，然后提起上面一块平板，由于分子力的作用，另一块平板将会粘在上面一块平板上并跟着被提起。这是一块平板上的原子对另一平板上的原子之间直接吸引的例证。

但是按照引力是基本的这个意义来说，这种分子吸引力还不算是基本的，它们是由一个分子中所有的电子和核与另一个分子中所有的电子和核之间的大量极其复杂的相互作用所引起的。我们得到的任何一个看起来简单的公式都是相当于复杂因素的总和，因此我们仍旧没有弄清楚基本的现象。

因为分子力在距离大时吸引，距离小时排斥，如图12.2所示，我们可以这样来形成固体，其中所有的原子依靠吸引力的作用结合在一起，而当原子靠得太近时，斥力就开始起作用使它们分开。在某一距离$d$处（图12.2中的曲线与轴相交的地方），作用力为零，这意味着所有的力都被平衡，因此分子与分子之间保持着这个距离。如果将分子推近到比距离$d$更近，分子就相互排斥，这就是$r$轴上方的曲线所表示的情况。要想把分子稍微推近一点就需要用很大的力，因为当距离小于$d$时，分子斥力迅速增大。如果分子被稍微拉开一点，就要有一点引力，引力随拉开的距离的增大而增加。如果分子被很大的力所拉开，它们将永远分开——键被拉断了。

如果分子相对于距离$d$仅被推近或拉开一段很小的距离，那么在图12.2曲线上相应移动的距离也是很小的，于是可以近似地用一条直线来表示。因此，在很多情况下，如果位移不是很大，力就与位移成正比。这条原理就是众所周知的胡克定律或弹性定律。此定律表明：当物体形变时，物体中试图恢复原状的力与形变的大小成正比。当然，此定律仅当形变较小时才是有效的。当形变太大时，物体将破裂或者被压碎，视形变的性质而定。为了使胡克定律成立，力的数值要有一定范围，它取决于材料的性质；例如，对面粉团或油灰来说，这个力是非常小的；但是对于钢，这个力就比较大。用一条垂直悬挂的钢制长螺旋弹簧可以很好地演示胡克定律。在弹簧的下端挂上适量的重物，可以使各处都产生很小的扭转，结果在每一匝中都引起一个小的垂直偏转，如果匝数很多，加起来就成为一个大的位移。比如说，如果测量100 g重物产生的总伸长，可以发现，每增加100 g重物将会产生一段附加伸长，它与相对于第一个100 g重物测得的伸长量几乎相等。当弹簧过载时，力与位移的这种定比关系开始变化，也就是说胡克定律不再有效。

\section{基本力、场}

下面我们来讨论唯一剩下的基本力。我们把它们称做基本力，是由于它们遵从的定律从根本上说是简单的。我们将首先讨论电力。物体带有仅由电子或质子组成的电荷。如果任何两个物体带上电荷，那么在它们之间就存在电力。如果电荷的大小分别是$q_1$和$q_2$，那么电力与两个电荷之间的距离的平方成反比，即$F=\frac{(\text{常数})q_1q_2}{r^2}$。对于异号电荷，这一定律与万有引力定律相似，但是，对于同号电荷，这个力是斥力，符号（方向）相反。本质上，电荷$q_1$和$q_2$可以是正的或负的。在公式的具体应用中，只要给予$q$以适当的正、负号就能得出力的正确的方向。力的方向沿着两个电荷的连线。当然，公式中的常数取决于力、电荷和距离所用的单位。通常，在实际应用中，电荷的单位是库仑（C），距离的单位是米（m），力的单位是牛顿（N）。为了使力恰好以牛顿为单位，常数（由于历史原因，这个常数被写成$\frac{1}{4 \pi \varepsilon_0}$）的值取
\begin{equation*}
    \varepsilon_0=8.85 \times 10^{-12} C^2 \cdot N^{-1} \cdot m^{-2}
\end{equation*}
或
\begin{equation*}
    \frac{1}{4 \pi \varepsilon_0}=8.99 \times 10^9 N \cdot m^2 \cdot C^{-2}
\end{equation*}
因此，静止电荷的作用力的定律为
\begin{equation}
    \label{Eq:I:12:2}
    \boldsymbol{F}=\frac{q_1q_2 \boldsymbol{r}}{4 \pi \varepsilon_0 r^3}
\end{equation}
在自然界中，所有电荷中最重要的是单个电子的电荷，它的电量为$1.60 \times 10^{-19}C$。在研究基本粒子之间的电力而不是研究大的电荷时，许多人宁可用$\frac{(q_{el})^2}{4π\varepsilon_0}$这种组合，其中$q_{el}$规定为电子的电荷。这种组合是经常出现的，为了简化计算，用记号$e^2$表示；在国际单位制中，它的数值是$(1.52 \times 10^{-14})^2$。采用这种形式常数的好处是，两个电子之间的力，用牛顿作为单位时可以简单地写成$\frac{e^2}{r^2}$；其中$r$的单位用米来表示，而不需要很多单独的常数。电力比这个简单公式所表示的要复杂得多，因为公式所给出的仅仅是当两个物体处于静止时，它们之间的力。接下去我们将讨论较普遍的情况。

在分析比较基本的一类力（不是像摩擦力，而是电力或引力之类的力）时形成了一种有趣的、非常重要的概念。因为乍看起来，力比反平方定律所指出的要复杂得多，而这些定律仅当相互作用物体处于静止时才成立，所以就需要一种改进的方法来处理当物体开始以一种复杂的方式运动时所产生的非常复杂的力。经验表明，用所谓“场”的概念这种方法，对于分析这种类型的力是非常有用的。比如说，以电力为例来说明这个概念，假定我们有两个电荷$q_1$、$q_2$分别位于$P$点和$R$点。那么两个电荷之间的力为
\begin{equation}
    \label{Eq:I:12:3}
    \boldsymbol{F}=\frac{q_1q_2 \boldsymbol{r}}{r^3}
\end{equation}
如果用场的概念来分析这个力，我们说$P$处的电荷$q_1$在$R$处产生了一种“状况”，而当电荷$q_2$被置于$R$处时，它就“感受”到这个作用力。这是一种描写力的方法，或许是奇特的。我们说，作用在$R$处电荷$q_2$上的力$\boldsymbol{F}$可以写成两部分。力等于电量$q_2$乘以一个量$\boldsymbol{E}$，不管有没有电荷$q_2$，$\boldsymbol{E}$都应当是存在的（只要我们将所有其他的电荷都保持在原来的位置上）。我们说$\boldsymbol{E}$是由$q_1$产生的“状况”，$\boldsymbol{F}$是$q_2$对于$\boldsymbol{E}$的响应。$\boldsymbol{E}$叫做电场，它是一个矢量。由$P$处电荷$q_1$在$R$处产生的电场$\boldsymbol{E}$的公式是电荷$q_1$乘以常数$\frac{1}{4 \pi \varepsilon_0}$除以$r^2$（$r$是$P$到$R$的距离），它作用于沿矢径的方向（矢径$\boldsymbol{r}$除以它自身的长度），因此$\boldsymbol{E}$的表示式为
\begin{equation}
    \label{Eq:I:12:4}
    \boldsymbol{E}=\frac{q_1 \boldsymbol{r}}{4 \pi \varepsilon_0 r^3}
\end{equation}
这样，我们就把力、电场和电场中电荷的关系式写成
\begin{equation}
    \label{Eq:I:12:5}
    \boldsymbol{F}=q_2 \boldsymbol{E}
\end{equation}
这样做的要点是什么？要点就是把分析分成两部分。一部分说，某物产生了一个场。另一部分说，某物受到场的作用。由于可以独立地看待这两部分，把分析分成这两部分在许多情况中简化了问题的计算。要是有许多电荷存在，我们可以先算出由所有电荷在$R$处产生的总电场，然后只要知道放在$R$处的电荷，我们就能求出作用在该电荷上的力。

对于引力的情形，我们完全可以同样处理。在这种情形下，力$\boldsymbol{F}=\frac{-Gm_1m_2\boldsymbol{r}}{r^3}$。我们可以作如下的类似分析：引力场中的物体所受的力等于物体的质量乘以引力场$\boldsymbol{C}$。$m_2$所受的力等于$m_2$的质量乘以由$m_1$所产生的场$\boldsymbol{C}$；即$\boldsymbol{F}=m_2\boldsymbol{C}$。质量为$m_1$的物体所产生的场$\boldsymbol{C}=\frac{-Gm_1 \boldsymbol{r}}{r^3}$，它的方向与电场的情况一样沿着径向。

不管初看起来如何，这种把一部分与另一部分分开的方法并不是微不足道的。如果力的定律真是简单的，它就没什么价值（只不过用另一种方法来写出同一件事情），但是力的定律是如此复杂，以至于结果表明场具有几乎与产生它们的物体无关的实在性。人们可以做某种事情，例如使电荷保持运动，于是在一定距离处产生一种效应——场；然而，如果电荷停止运动，场记录着过去所有的情况，因为两个粒子之间的相互作用不是瞬时的。我们希望有某种方法记住以前所发生的事情。如果某电荷所受的力取决于另一个电荷昨天所在的位置以及它当时的行为，那么我们就需要一种机构来记录昨天发生的事情，这就是场的特征。因此，当力变得越复杂，场就变得越来越真实，而这种分离技巧的人为性也就越来越少。

用场来分析力时，我们需要用到有关场的两种定律。第一种是对场的响应，它给出了运动方程。例如，质量对引力场的响应定律为：力等于质量乘以引力场；或者，如果物体还带有电荷，电荷对于电场的响应等于电荷乘以电场。对这些情况的性质的第二部分分析是把场的强度以及它是如何产生的规律用公式表示出来。有时，把这些定律称为场方程。在适当的时候，我们将更深入地学习场方程，这里我们只写出几点有关的内容。

第一，一切事实中最为惊人的是，由若干源产生的总电场是由第一个源、第二个源等等产生的电场的矢量和，这是完全确实而又容易理解的。换句话说，如果有许多电荷产生一个场，则其中的一个电荷独自产生的电场为$\boldsymbol{E}_1$，另一个电荷独自产生的电场为$\boldsymbol{E}_2$等等。那么，只要把所有的矢量加起来就得到了总电场。这个原理可以表示成
\begin{equation}
    \label{Eq:I:12:6}
    \boldsymbol{E}=\boldsymbol{E}_1+\boldsymbol{E}_2+\boldsymbol{E}_3+\cdots
\end{equation}
或者，根据上面给出的定义
\begin{equation}
    \label{Eq:I:12:7}
    \boldsymbol{E}=\sum_{i} \frac{q_ir_i}{4\pi \varepsilon_0 r^3_i}
\end{equation}

同样的方法适用于引力吗？牛顿把两个物体$m_1$和$m_2$之间的力表示成$\boldsymbol{F}=\frac{-Gm_1m_2 \cdot \boldsymbol{r}}{r^3}$。但是按照场的概念，我们可以说物体m1在其周围空间产生了引力场$\boldsymbol{C}$，$m_2$所受的力则由
\begin{equation}
    \label{Eq:I:12:8}
    \boldsymbol{E}=m_2\boldsymbol{C}
\end{equation}
给出。同电场的情况完全类似
\begin{equation}
    \label{Eq:I:12:9}
    \boldsymbol{C}_i=\frac{-Gm_i\boldsymbol{r}_i}{r^3_i}
\end{equation}
由几个物体产生的引力场为
\begin{equation}
    \label{Eq:I:12:10}
    \boldsymbol{C}=\boldsymbol{C}_1+\boldsymbol{C}_2+\boldsymbol{C}_3+\cdots
\end{equation}
在第9章中计算行星运动情况时，实质上我们已经应用了这一原理。我们正是把所有的力矢量加起来以得到作用在一个行星上的合力。要是约去该方程中行星的质量，我们就得到式（12.10）。

式（12.6）和（12.10）表示的就是所谓场的叠加原理。这个原理说明，由所有的源产生的总的场等于由每一个源产生的场之和。就我们目前所知，对于电学这是一个绝对保证的定律。甚至由于电荷运动而使力的定律变为复杂时，这个定律仍然正确。有一些表面上的反例，但只要再仔细分析一下，总会发现这是由于忽略了某些运动电荷。然而，虽然叠加原理完全适用于电力，但是对于很强的引力场，叠加原理不完全正确。按照爱因斯坦的引力场理论，牛顿方程式（12.10）仅是一种近似。

与电力紧密相关的另一类力称为磁力，这种力也是用场来进行分析的。电力和磁力之间的一些定性关系可以用一个电子射线管（图12.3）的实验来说明。电子射线管的一端是
\begin{wrapfigure}{l}{0.4\textwidth}
    \centering
    \includegraphics[width=0.35\textwidth]{Chapter12/电子束管}
    \caption{电子束管}
    \label{figure:电子束管}
\end{wrapfigure}
一个发射电子流的源。在管子里面有一套装置把电子加速到很高的速度，并聚焦成很窄的电子束再送到管子另一端的荧光屏上。在荧光屏的中央，电子打到的地方发出一个亮的光点，这样我们就能够跟踪电子的径迹。在射向荧光屏的途中，电子束穿过一对水平放置的平行金属板中间的窄缝。两块金属板上可以加上电压，因此可以随意成为带负电的。当加上电压时，两块金属板之间就产生一个电场。

实验的第一部分是给下极板加上负电压，这意味着把额外的电子放到下极板上。由于同种电荷相互排斥，荧光屏上的光点立即向上移动（我们也可以用另一种方式来表明——电子“感受”到场，并以向上偏转作为响应）。接着，我们把电压极性反过来，使上极板为负。现在，荧光屏上的光点跳到中心之下，这表明电子束中的电子受到上极板中电子的排斥（或者我们又可以说电子对场作出响应，现在是在相反的方向上）。

实验的第二部分是切断极板上的电压，试验磁场对电子束的影响。这一步是用一个马蹄形磁铁来进行的，磁铁的两极分得足够开，可以或多或少地跨立在管子上。假定我们把磁铁以字母$U$那样的取向放到管子下面，两极向上，使管子的一部分位于磁极之间。我们看到光点偏转了，比如说，当磁铁从下方趋近管子时，光点就向上偏转。这样看来好像是磁铁排斥电子束。然而事情并不那样简单，如果我们把整个磁铁倒转过来，但磁极的位置并不对调，那么从上方趋近管子，光点还是向上移动。由此看来，电子束不是受到排斥；相反，好像是受到吸引。现在，我们再从头开始，把磁铁恢复到原来的$U$形取向，并把它放在管子下方。没有错，光点还是向上偏转。现在将磁铁绕垂直轴旋转180°，这样，磁铁仍处于$U$形位置，但两个磁极的相对位置颠倒了。瞧！现在光点跳向下方，并停在下方，即使像前面那样把磁铁倒转过来，从上方趋近管子，光点还是停在下方。

要理解这种独特的行为，我们必须有一种新的力的组合，因此，我们把它解释为：在磁铁的两极之间存在着磁场。这种场是有方向性的，其方向总是由一特定的极（我们可以标上记号）出发，指向另一极。将磁铁倒转过来并不改变场的方向，但是将一磁极相对另一磁极颠倒一下就使场的方向变得相反。举例来说，如果电子速度沿$x$方向是水平的，则磁场也是水平的，但在$y$方向上，作用于运动电子上的磁力应在$z$方向上，即向上或向下，取决于磁场是在正$y$方向还是负$y$方向。

尽管目前我们不准备介绍彼此以任意方式相对运动的带电物体之间力的正确定律，但因为定律太复杂，我们将介绍它的一个方面：如果场是已知时力的完整定律。作用在带电物体上的力取决于它的运动；如果物体在一给定位置上停止不动时，有某个力作用着，则这个力与电荷量成正比，比例系数就是我们所谓的电场。当物体运动时，力就不同了，我们发现其修正项，即新的“一份”力严格线性地取决于速度，并与$v$和另一个我们称为磁感应强度$\boldsymbol{B}$的矢量正交。如果电场$\boldsymbol{E}$和磁感应强度$\boldsymbol{B}$的分量分别为$(E_x,~E_y,~E_z)$和$(B_x,~B_y,~B_z)$，速度$v$的分量为$(v_x,~v_y,~v_z)$，那么作用在运动电荷$q$上总的电磁力的分量为
\begin{equation}
    \label{Eq:I:12:11}
    \begin{split}
        F_x=q(E_x+v_yB_z-v_zB_y) \\
        F_y=q(E_y+v_zB_x-v_xB_z) \\
        F_z=q(E_z+v_xB_y-v_yB_x)
    \end{split}
\end{equation}
举例来说，如果仅有磁场分量$B_y$和速度分量$v_x$，那么磁力项中留下的只有$z$方向的力，它与$\boldsymbol{B}$和$v$两者正交。

\section{赝力}

下面，我们将讨论的这一类力可以称为赝力。在第11章中，我们讨论了使用不同坐标系的乔和莫两个人之间的关系。我们假定，粒子的位置由乔测得的是$x$，由莫测得的是$x'$；于是定律如下所示
\begin{equation*}
    x=x'+s,~y=y',~z=z'
\end{equation*}
式中$s$是莫的坐标系相对于乔坐标系的位移。如果我们假定，运动定律对于乔是正确的，那么定律在莫看来又如何呢？首先，我们发现
\begin{equation*}
    \frac{dx}{dt}=\frac{dx'}{dt}+\frac{ds}{dt}
\end{equation*}
前面，我们考虑了$s$是常量的情况，我们发现$s$对运动定律毫无影响，因为$\frac{ds}{dt}=0$；因此，最终物理定律在两种坐标系中是相同的。但是我们可以取的另一种情况是$s=ut$，式中$u$是沿直线运动的均匀速度，于是$s$不是常量，$\frac{ds}{dt}$就不为零，但等于一个常数$u$。然而，加速度$\frac{d^2x}{dt^2}$仍然与$\frac{d^2x'}{dt^2}$相同，因为$\frac{du}{dt}=0$。这就验证了我们在第10章中所使用的定律，即如果我们沿直线以均匀速度运动，则我们所看到的物理定律与处于静止时相同。这就是伽利略变换。但是，我们希望讨论$s$的更加复杂的有趣情况，比如说$s=\frac{at^2}{2}$。于是$\frac{ds}{dt}=at$，而$\frac{d^2s}{dt^2}=a$，一个均匀加速度；或者更复杂一些，加速度可以是时间的函数。这就意味着，虽然从乔的观点来看运动定律应为
\begin{equation*}
    m\frac{d^2x}{dt^2}=F_x
\end{equation*}
而从莫看来则应是
\begin{equation*}
    m\frac{d^2x'}{dt^2}=F_{x'}-ma
\end{equation*}
也就是说，由于莫的坐标系相对于乔的坐标系是加速的，因此出现了额外项$ma$，而莫就必须用这个量来修正他的力，以便使牛顿定律继续有效。换句话说，这是一个明显的、起源不明的神秘的新力，当然它的出现是因为莫使用了不正确的坐标系。这是赝力的一个例子；其他的例子出现在转动的坐标系中。

赝力的另一个例子是通常所谓的“离心力”。在转动坐标系中，例如一个旋转的箱子里的观察者在把东西抛向墙壁时，将会发现一种神秘的力，这个力用任何已知起源的力都解释不了。这些力的出现，只是由于观察者不具备牛顿坐标系这个事实，而牛顿坐标系是最简单的坐标系。

赝力可以通过一个有趣的实验来说明，在这个实验中，我们拖一桶水沿着一张桌子加
\begin{wrapfigure}{l}{0.4\textwidth}
    \centering
    \includegraphics[width=0.35\textwidth]{Chapter12/赝力的实例}
    \caption{赝力的实例}
    \label{figure:赝力的实例}
\end{wrapfigure}
速前进。当然，作用在水上的重力方向向下，但是由于水平方向的加速，也有一个水平作用着的赝力，方向与加速度方向相反。重力与赝力的合力与垂直方向成一角度，在加速运动期间，水的表面将与合力垂直，即与桌面倾斜成一角度，在桶的后面部分水面将高起一些。当我们不再拖水桶时，由于摩擦而使桶减速运动（赝力变换了方向），桶前面部分的水将高起一些（图12.4）。

赝力的一个非常重要的特征是，它们永远与质量成正比；重力也是这样。因此，有可能重力本身就是一种赝力。或许简单地说，引力就是由于我们没有正确的坐标系而引起的，这难道是不可能的吗？归根到底，如果我们设想一个物体正在加速，那么我们总可以得到一个与质量成正比的力。例如，一个关闭在箱子里的人（箱子静止放在地球上），就会发现有一个力使他待在箱子的地板上，这个力与他的质量成正比。但是如果根本没有地球，而箱子静止不动，那么箱子里的人就会漂浮在空中。另一方面，如果根本没有地球，而是有某个东西将箱子以加速度$g$向上拉，那么箱子里的人在分析物理现象时，应当发现一个赝力，这个力就像重力那样把人拉向地板。

爱因斯坦提出了著名的假设，加速度产生引力的赝物，加速度的力（赝力）与引力是不可能区分的；要说出给定的力中有多少是重力，有多少是赝力是不可能的。

把重力看作赝力，比如说我们都保持向下是由于我们在向上加速，这似乎没有问题，但是在地球另一边的马达加斯加人会怎样呢？——他们也在加速吗？爱因斯坦发现，每次只有在一个点上才可以把重力同时看成赝力，根据这个考虑，他认为世界的几何性要比普通欧几里得几何复杂得多。我们现在的讨论仅仅是定性的，并不想涉及超出一般概念之外的东西。为了对引力怎么会是赝力的结果有个大致的概念，我们来作一个纯粹是几何的，并不代表真实情况的说明。假定，我们大家都在二维空间中生活，对第三维空间毫无所知。我们只认为，我们处在一个平面上，但是假定我们实际上处在一个球的表面上。再假定我们沿地面发射出一个物体，物体上没有力作用着。它将到哪儿去呢？看来它似乎沿直线运动，但是必须仍处在球的表面上，球面上两点之间的最短距离是沿着一个大圆的；于是它就沿大圆运动。如果我们同样地发射另一个物体，但沿另一方向，它就沿另一大圆运动。因为我们认为，我们是处在一个平面上，可以预期，这两个物体之间的距离将随时间线性地增加，但经过仔细观察将会发现，如果两个物体运动得足够远，那么它们又会互相接近，仿佛相互吸引。但是它们并不相互吸引——关于这种几何学真是有点“古怪”。这个特定的例子并没有正确地描述爱因斯坦几何学“古怪”在哪里，但是它说明，如果我们使几何形状充分畸变，那么所有的引力以某种方式与赝力联系起来是可能的；这就是爱因斯坦引力理论的一般观念。

\section{核力}

作为本章的结束，我们简短地讨论一下仅有的另外一些已知的力，即所谓的核力。这些力存在于原子核的内部，尽管对它们进行了充分的讨论，但没有一个人曾经计算过两个原子核之间的力，甚至在目前还没有关于核力的已知定律。这些力的作用范围极小，差不多与原子核的大小（大约是$10^{-13}cm$）相同。对于这样小的粒子，再加上作用距离又是如此微小，只有量子力学的定律才是正确的，牛顿定律就失效了。在核分析中，我们不再用力来思考，事实上我们可以用两个粒子的相互作用能的概念来代替力的概念，这一课题将在以后进行讨论。关于核力，能够写出的任何公式都是忽略了许多复杂情况的相当粗糙的近似；其中之一可以表述如下：在原子核之内的力并不随距离的平方反比变化，而是在一定距离$r$之后指数地衰减掉，用公式表示，即$F=(\frac{1}{r^2})exp(-\frac{r}{r_0})$，式中距离$r_0$的数量级是$10^{-13}cm$。换句话说，粒子相隔距离稍大一些，力就消失了，虽然这些力在$10^{-13}cm$范围内是非常强的。就今天对核力的理解而言，它的定律是非常复杂的；我们不能以简单的方式去理解它们，因而分析核力背后的基本机理的整个问题仍未解决。在试图解决这个问题的过程中，发现了许多奇异粒子，例如$\pi$介子，但是这类力的起因仍然不清楚。